%% kapa: exam=1 round=28 year=2017 period=2 form=A subject=appraisal_law subject_ko=감정평가관계법규 seq=14 number=14
%% source: Q-Net 감정평가사 기출문제 (cst00309) artlSeq=5199603 file=extracted/1008/5199603/2017년 제28회 감정평가사 1차시험 문제지/2교시 A형.pdf
%% conversion: kordoc 4.24.51 -> tools/gen1.py (auto); review=proofread-sonnet; flags=none
\begin{question}{14}
감정평가 및 감정평가사에 관한 법령상 감정평가에 관한 설명으로 옳지 않은 것은?
\begin{choices}
  \item 감정평가업자는 해산하거나 폐업하는 경우 감정평가서의 원본과 그 관련 서류를 국토교통부장관에게 제출하여야 한다.
  \item 감정평가업자는 감정평가서의 관련 서류를 발급일로부터 5년 이상 보존하여야 한다.
  \item 감정평가 의뢰인이 감정평가서를 분실하거나 훼손하여 감정평가서 재발급을 신청한 경우 감정평가업자는 정당한 사유가 있을 때를 제외하고는 감정평가서를 재발급 하여야 한다.
  \item 국가가 토지등을 경매하기 위하여 감정평가를 의뢰하려고 한국감정평가사협회에 감정평가업자 추천을 요청한 경우 협회는 요청을 받은 날부터 7일 이내에 감정평가업자를 추천하여야 한다.
  \item 유가증권도 감정평가의 대상이 된다.
\end{choices}
\end{question}
\answer{2}
